\chapter{Accéder aux marchés~: actions}\label{ch:m1:getting-access-equities}

Trois personnes disposent d'une stratégie qui fonctionne sur le papier, de dix millions de dollars
de leurs propres fonds et de ceux de leurs investisseurs, et n'ont aucun moyen d'envoyer un ordre.
Aucune bourse n'acceptera un ordre de leur part~: les bourses traitent avec des membres.
Aucun membre ne prêtera sa qualité de membre sans contrôles, car la loi rend le membre
responsable de chaque ordre qui transite par lui. Personne ne réglera leurs transactions
tant qu'une société de compensation n'aura pas accepté de les garantir.
Les données dont elles ont besoin pour décider de ce qu'il faut négocier sont soumises à licence,
et la licence d'une machine qui lit ces données coûte plus cher que celle d'une personne.
Et le prix de chaque action qu'elles négocient, au centième de cent près, dépendra
du nombre d'actions échangées dans le mois. Les vingt-huit premiers chapitres ont décrit le marché
comme un mécanisme. Ce chapitre et le suivant le décrivent comme un ensemble de contrats,
ce qui est la manière dont une société le rencontre.

\section{Membre ou client~?}

\begin{definition}[Courtiers exécutant et compensateur]\label{def:m1:getting-access-equities:brokers}
Un \emph{courtier exécutant}\index{courtier exécutant} est le membre par l'intermédiaire duquel
les ordres d'une société atteignent une plateforme, au nom du courtier et sous
ses contrôles de risque pré-négociation (\cref{ch:m1:exchanges-brokers-venues}). Un \emph{courtier
compensateur}\index{courtier compensateur} est le membre de la chambre de compensation qui règle
les transactions de la société, porte ses positions et les finance. Un
\emph{give-up}\index{give-up (equities)} est une transaction exécutée par un courtier puis
transmise, par accord, à un autre pour compensation, afin qu'une société puisse
exécuter auprès de plusieurs courtiers et tout détenir chez un seul.
\end{definition}

Il existe trois voies, par ordre croissant de coût, de contrôle et de rapidité.

\textbf{Client d'un courtier.} La société envoie ses ordres aux algorithmes ou
au routeur intelligent d'un courtier. Aucune inscription~; ce sont les frais, les paliers et la latence du courtier.

\textbf{\omterm{def:m1:exchanges-brokers-venues:dma}{Accès direct au marché} sponsorisé.} Les ordres de la société parviennent à la
bourse par l'intermédiaire de l'adhésion du courtier et, en vertu de la loi, par des
contrôles de risque que le courtier maîtrise~; la société choisit elle-même la plateforme et le type d'ordre, souvent
depuis ses propres serveurs installés dans le centre de données de la bourse. Elle paie le courtier par
action et bénéficie du palier de bourse accordé au courtier, en totalité ou en partie.

\textbf{Inscription et adhésion.} La société devient un \omterm{def:m1:the-sell-side:sell-side}{broker-dealer} et un
membre de chaque bourse dont elle a besoin. Elle fait directement face à la bourse, perçoit ses propres
paliers, peut rejoindre des programmes de fourniture de liquidité et assume le capital,
la déclaration, la supervision et le contrôle.

\begin{dated}{2026-09}{Ce que coûte l'inscription en capital}\label{dat:m1:getting-access-equities:netcap}
Selon la règle du capital net, un \omterm{def:m1:the-sell-side:sell-side}{broker-dealer} qui tient des comptes clients
doit maintenir un capital net d'au moins 250\,000~\$~; un dealer, catégorie qui comprend toute
société réalisant plus de dix transactions par an pour son propre compte d'investissement, au moins
100\,000~\$~; un courtier qui présente ses clients à une société de compensation,
50\,000~\$~; celui qui ne détient jamais de fonds ni de titres de clients, 5\,000~\$. Un
\omterm{def:m1:what-a-trading-firm-does:market-maker}{teneur de marché} a besoin de 2\,500~\$ pour chaque titre pour lequel il assure la tenue de marché
(1\,000~\$ pour ceux cotés à 5~\$ ou moins), dans la limite de 1~million de dollars. L'endettement agrégé
ne peut dépasser 1\,500~\% du capital net. Il s'agit de planchers~; la
contrainte effective pour une société de trading est la décote appliquée à ses positions.
\end{dated}

\begin{omfigure}
\begin{tikzpicture}[font=\small,
  box/.style={draw=omDef,rounded corners=2pt,minimum width=2.6cm,minimum height=0.95cm,align=center,fill=omDef!4},
  ext/.style={draw=omExo,rounded corners=2pt,minimum width=2.6cm,minimum height=0.95cm,align=center,fill=omExo!6},
  flow/.style={-{Stealth[length=2mm]},thick}]
\node[box] (f) at (0,0) {Société de trading};
\node[box] (eb) at (4.3,1.3) {Courtier exécutant\\(contrôles de risque)};
\node[ext] (ex) at (9.8,1.3) {Bourses};
\node[box] (cb) at (4.3,-1.3) {Courtier compensateur\\(financement, marge)};
\node[ext] (ccp) at (9.8,-1.3) {Chambre de compensation};
\draw[flow,omDef] (f.north east) -- node[above,sloped,font=\footnotesize]{ordres} (eb.west);
\draw[flow,omDef] (eb) -- node[above,font=\footnotesize]{comme membre} (ex);
\draw[flow,omThm] (ex) -- node[right,font=\footnotesize]{transactions} (ccp);
\draw[flow,omThm] (ccp) -- node[above,font=\footnotesize]{règle} (cb);
\draw[flow,omThm] (cb.west) -- node[below left,font=\footnotesize,align=right,pos=0.5]{positions,\\espèces} (f.south east);
\draw[flow,omExo,dashed] (eb) -- node[right,font=\footnotesize]{give-up} (cb);
\end{tikzpicture}

\omcaption{Une société qui n'est ni membre ni membre compensateur. Les ordres remontent
par le \omterm{def:m1:getting-access-equities:brokers}{courtier exécutant}~; les transactions redescendent par le
\omterm{def:m1:getting-access-equities:brokers}{courtier compensateur}, qui est la véritable contrepartie de la société pour les espèces
et les titres.}\label{fig:m1:getting-access-equities:chain}
\end{omfigure}

\section{La convention de compensation}

Le \omterm{def:m1:getting-access-equities:brokers}{courtier compensateur} est le contrat le plus important, et celui auquel les nouvelles sociétés
pensent en dernier. Il détermine le levier dont dispose la société (sa marge maison, et non
le minimum réglementaire), ce qu'elle paie pour financer les positions acheteuses et emprunter
les positions vendeuses (\cref{ch:m1:financing}, \cref{ch:m1:stock-loan-and-short-selling}),
si les positions sont margées en portefeuille, ce qu'il facture par action ou par ticket
pour compenser, et les conditions dans lesquelles il peut relever la marge ou liquider
sans préavis. Un \omterm{def:m1:getting-access-equities:brokers}{courtier compensateur} prend le risque de crédit de la société pendant la nuit~;
il voudra examiner le profil de risque de la stratégie, l'historique des principaux,
le capital audité et un engagement de revenu minimal. Une société dont le \omterm{def:m1:getting-access-equities:brokers}{courtier compensateur}
accorde un préavis de trente jours dispose de trente jours pour en trouver un autre
ou pour cesser son activité.

\section{Grilles tarifaires et paliers}

\begin{definition}[Palier de volume]\label{def:m1:getting-access-equities:tier}
Un \emph{palier de volume}\index{palier de volume} est une ligne de la grille tarifaire d'une bourse
qui accorde un meilleur tarif aux membres dont l'activité mensuelle dépasse un seuil,
exprimé en général comme le volume moyen quotidien ajouté rapporté au volume consolidé total
du marché. Le tarif du palier atteint
s'applique à l'ensemble des actions éligibles du mois.
\end{definition}

\begin{dated}{2026-09}{La grille d'une bourse}\label{dat:m1:getting-access-equities:bzx}
La grille tarifaire des actions d'une bourse américaine à modèle maker-taker, en vigueur le 1\up{er}~septembre 2026,
pour les actions cotées 1~\$ ou plus~: le retrait de liquidité coûte 0.0030~\$ par action~; l'ajout de liquidité affichée
rapporte 0.0016~\$ au tarif standard et, dans le cadre des paliers d'ajout de volume, 0.0020~\$ à partir de 0.06~\%
du volume consolidé total, 0.0023~\$ à partir de 0.20~\%, 0.0028~\$ à partir de 0.25~\% et, au palier supérieur
mentionné ici, 0.0031~\$ à partir de 1.00~\%. Les ordres non affichés rapportent de zéro
à 0.0008~\$. La grille comporte de nombreux autres paliers, par bande de cotation,
par type d'ordre et pour des combinaisons d'activité sur
les bourses du groupe.
\end{dated}

\begin{omfigure}
\begin{tikzpicture}
\begin{axis}[width=11cm,height=5cm,
  xlabel={volume moyen ajouté par jour (millions d'actions)},ylabel={milliers de \$ par mois},
  tick label style={font=\small},label style={font=\small},
  xmin=0,xmax=40,ymin=0,ymax=2500,grid=major,grid style={gray!25},
  scaled y ticks=false,yticklabel style={/pgf/number format/fixed,/pgf/number format/1000 sep={\,}},
  table/col sep=comma]
\addplot[omDef,very thick,const plot] table[x=adav_m,y=monthly_k]{figdata/markets-1/29-getting-access-equities/tiers.csv};
\node[font=\footnotesize,anchor=west] at (axis cs:7.6,230) {0.06~\%};
\node[font=\footnotesize,anchor=west] at (axis cs:24.4,1030) {0.20~\%};
\node[font=\footnotesize,anchor=west] at (axis cs:30.4,1560) {0.25~\%};
\end{axis}
\end{tikzpicture}

\omcaption{Ristournes mensuelles selon la grille de
l'\cref{dat:m1:getting-access-equities:bzx}, avec un volume consolidé total
supposé de 12 milliards d'actions par jour et 21 jours de négociation. Le taux d'un palier
s'appliquant à l'ensemble du mois, chaque seuil est une falaise~: l'action qui franchit
le seuil de 0.20~\% vaut 155\,000~\$. Données~: la construction du chapitre.}\label{fig:m1:getting-access-equities:tiers}
\end{omfigure}

Trois conséquences en découlent. \emph{La taille est un avantage}~: deux sociétés dotées de la même
stratégie gagnent des montants différents par action, et la plus grande peut coter là où la
plus petite ne le peut pas. \emph{L'agrégation est un métier}~: un courtier dont les clients
atteignent ensemble un palier élevé peut reverser l'essentiel de ce taux à chacun d'eux et
garder le reste, ce qui est souvent une meilleure affaire pour une petite société que son propre
statut de membre. \emph{Un volume proche d'un seuil a un prix fictif}~: en fin
de mois, une société juste en deçà d'un palier négociera des actions à perte pour l'atteindre,
et son comportement change le jour où elle l'atteint.

\begin{proposition}[Quand le remplissage de volume est rentable]\label{prop:m1:getting-access-equities:padding}
Un membre ajoute naturellement $v$ actions par jour, en touchant une ristourne $\rho_0$ par action.
Le palier suivant, à $V^* > v$, verse $\rho_1 > \rho_0$. Du volume supplémentaire peut être obtenu
à une perte de $\ell$ par action avant ristournes, avec $\ell > \rho_1$. Le remplissage jusqu'à $V^*$
est rentable si et seulement si
\[
v \;>\; V^*\,\frac{\ell - \rho_1}{\ell - \rho_0}.
\]
\end{proposition}
\begin{proof}
Le remplissage modifie le résultat quotidien de $V^*\rho_1 - (V^*-v)\ell - v\rho_0$, qui
est positif exactement lorsque $v(\ell-\rho_0) > V^*(\ell-\rho_1)$.
\end{proof}

\begin{omfigure}
\begin{tikzpicture}
\begin{axis}[width=10.5cm,height=4.8cm,ybar,bar width=12pt,
  ylabel={cents pour 100 actions},
  symbolic x coords={maker-taker,inverted,flat-fee},xtick={maker-taker,inverted,flat-fee},xticklabels={{maker-taker},{inversée},{frais fixes}},
  tick label style={font=\small},label style={font=\small},xticklabel style={font=\small\strut},
  ymin=-25,ymax=42,ymajorgrids,grid style={gray!25},enlarge x limits=0.22,
  nodes near coords,nodes near coords style={font=\footnotesize,fill=white,inner sep=1pt,/pgf/number format/.cd,fixed,precision=0},
  legend columns=2,legend style={font=\footnotesize,at={(0.5,-0.2)},anchor=north,draw=none,fill=none,/tikz/every even column/.append style={column sep=8pt}},
  table/col sep=comma]
\addplot[fill=omDef!60,draw=omDef] table[x=venue,y=passive_mils]{figdata/markets-1/29-getting-access-equities/venues.csv};
\addplot[fill=omThm!60,draw=omThm] table[x=venue,y=aggressive_mils]{figdata/markets-1/29-getting-access-equities/venues.csv};
\legend{ordre au repos exécuté,ordre marketable}
\end{axis}
\end{tikzpicture}

\omcaption{Coût total par 100 actions sur trois types de plateformes pour un membre
ajoutant 9 millions d'actions par jour~: frais ou ristourne de la bourse, plus 2 cents de
compensation et 1 cent de frais réglementaires. Un montant négatif est un revenu. Les taux
maker-taker sont ceux publiés~; la \omterm{def:m1:us-equity-market-structure:makertaker}{plateforme inversée} et la plateforme à frais fixes, ainsi que les deux autres
coûts sont illustratifs. Données~: le tutoriel.}\label{fig:m1:getting-access-equities:venues}
\end{omfigure}

Le signe des frais détermine où un ordre doit reposer. Sur la \omterm{def:m1:us-equity-market-structure:makertaker}{plateforme inversée}, un ordre au repos
\emph{paie}, de sorte que la file est courte et que les exécutions arrivent
en premier~; sur la plateforme maker-taker, il est rémunéré, donc la file est longue et l'ordre est
exécuté en dernier, au moment où le prix est sur le point de le traverser. La différence de frais,
trois dixièmes de cent, est le prix que le marché facture pour la priorité
dans la file, et un routeur qui l'ignore
(\cref{ch:m1:us-equity-market-structure}) sous-évalue chaque ordre passif.

\section{Programmes de fournisseurs de liquidité}

\begin{definition}[Teneur de marché désigné et fournisseur de liquidité supplémentaire]\label{def:m1:getting-access-equities:dmm}
Un \emph{teneur de marché désigné}\index{teneur de marché désigné} (designated market maker) est le
membre assigné à un titre coté, avec l'obligation d'y maintenir un marché équitable et
ordonné, y compris lors des enchères d'ouverture et de clôture, en contrepartie de
privilèges d'information et d'allocation. Un \emph{fournisseur de liquidité
supplémentaire}\index{fournisseur de liquidité supplémentaire} (supplemental liquidity provider) est un membre qui s'engage à coter
de manière compétitive les titres qui lui sont assignés pendant une part minimale de la journée, en contrepartie de
ristournes majorées, sans les autres
obligations du \omterm{def:m1:what-a-trading-firm-does:market-maker}{teneur de marché} désigné.
\end{definition}

\begin{dated}{2026-09}{Une obligation de cotation}\label{dat:m1:getting-access-equities:slp}
À la New York Stock Exchange, un \omterm{def:m1:getting-access-equities:dmm}{fournisseur de liquidité supplémentaire} doit maintenir un cours
acheteur ou un cours vendeur au meilleur cours acheteur ou au meilleur cours vendeur national sur chaque
titre assigné pendant au moins 10~\% de la journée de négociation~; il peut s'agir d'une unité de
trading pour compte propre d'un membre ou d'un \omterm{def:m1:what-a-trading-firm-does:market-maker}{teneur de marché} agréé, et les deux ne sont pas
agrégés pour cette exigence.
\end{dated}

Des programmes de ce type existent sur chaque bourse, sous des noms différents.
Ce sont des contrats~: des obligations mesurées par la bourse (temps passé au meilleur cours, taille
cotée, fourchette), en échange de taux inaccessibles par ailleurs. Une société qui cote déjà
au meilleur cours la plus grande partie de la journée doit faire partie du programme~; une société
qui devrait changer de comportement pour se qualifier vend une option sur ses cotations contre la
ristourne, et doit la valoriser.

\section{Licences de données}

\begin{definition}[Redevance non-display]\label{def:m1:getting-access-equities:nondisplay}
Une \emph{redevance non-display}\index{redevance non-display} est le montant qu'une bourse
facture pour l'utilisation de ses données de marché par une machine (un algorithme, un routeur, un
système de risque), par opposition à une personne qui regarde un écran. Elle est prélevée
par société ou par plateforme, par catégorie d'usage, en sus des frais de réception des
données.
\end{definition}

La facture de données d'une société de trading comporte quatre niveaux~: l'accès
(le flux lui-même et les ports et interconnexions physiques par lesquels il arrive), la
redistribution (si les données vont à d'autres entités, même affiliées), l'affichage (par utilisateur) et le
non-display (par catégorie d'usage automatisé). Les bourses procèdent à des audits. La facture est établie par
bourse, et une stratégie qui a besoin des \omterm{def:m1:us-equity-market-structure:sip}{flux directs} d'une douzaine de bourses
(\cref{ch:m1:reading-market-data}) en paie une douzaine avant sa première transaction. Les données de marché représentent
une part importante des revenus d'un groupe boursier
(\cref{ch:m1:exchanges-brokers-venues}) et, pour une petite société, de ses coûts.

\section{La négociation}

Rien de ce qui figure dans un barème publié d'une bourse n'est négociable~: les frais des bourses sont déposés auprès du
régulateur et s'appliquent à chaque membre qui remplit leurs conditions,
ce qui explique que les barèmes soient longs. Tout le reste l'est~: la rétrocession à la société de son palier par le courtier~; les frais de
compensation par action et leur minimum~; les \omterm{def:m1:delta-one-instruments:swap}{spreads de financement} et les taux d'emprunt~;
la part des revenus de \omterm{def:m1:financing:seclending}{prêt de titres} du \omterm{def:m1:getting-access-equities:brokers}{courtier compensateur} sur les positions longues de la société
qui lui revient~; les prix des interconnexions et de l'hébergement dans les
centres de données de tiers~; les données et logiciels des fournisseurs. Ce que la société apporte à chaque
négociation, c'est son volume, ses soldes et sa capacité crédible d'aller ailleurs. Elle
doit donc connaître ses propres chiffres mieux que sa contrepartie~: nombre d'actions par jour par plateforme
et par type d'ordre (ajout ou retrait)~; soldes moyens, longs et
courts~; usage des titres difficiles à emprunter~; messages par seconde. Les constructions de ce livre
en produisent la plupart.

\section{Tutoriel~: coût net par action sur trois plateformes}\label{tut:m1:getting-access-equities:cost}

\begin{tutorial}
\textbf{Objectif.} Encoder un barème à paliers, trouver le palier d'un membre, calculer la facture du
mois, mesurer la rupture de coût à un seuil, et comparer le coût tout compris d'une action passive
et d'une action agressive sur trois types de plateformes. \textbf{État
final~:} les deux chiffres de données de ce chapitre.
\begin{tutsteps}
\item \textbf{Un barème} est un taux de base, un taux de retrait et une liste de paliers.
\omcode{code/firm/feesched/firm_feesched.py}{18}{41}{Recherche du palier et facture quotidienne.}\label{lst:m1:getting-access-equities:schedule}
\item \textbf{Remplissage.} Le meilleur taux s'applique à tout~; les actions de
remplissage perdent de l'argent avant la ristourne.
\omcode{code/firm/feesched/firm_feesched.py}{44}{54}{Vaut-il la peine d'ajouter du volume pour atteindre le palier suivant~?}\label{lst:m1:getting-access-equities:padding}
\item \textbf{Tout compris.} Frais de bourse, de compensation et réglementaires.
\end{tutsteps}
\textbf{Pour aller plus loin.} Ajouter une seconde condition à un palier (une part minimale
de volume \emph{retiré} également) et un palier inter-bourses qui prend en compte
l'activité sur une bourse sœur~; puis trouver l'allocation la moins chère d'un flux
d'ordres fixe entre les bourses du groupe.
\end{tutorial}

\section{Construction~: le moteur de barèmes de frais}\label{bld:m1:getting-access-equities:feesched}

\begin{build}
\bfield{Rôle} Le backtester de la société miniature (livre~7) et
le routeur (livre~11) ont besoin des frais marginaux de chaque ordre, qui dépendent du
volume du mois écoulé~; sa direction financière a besoin de la facture.
\bfield{Interface} \texttt{Tier(name, min\_adav\_share, add\_rate)};
\texttt{Schedule(venue, base\_add\_rate, remove\_rate, tiers)} avec
\texttt{tier\_for}, \texttt{add\_rate}, \texttt{daily\_bill}, \texttt{next\_tier}; \texttt{extra\_volume\_worth\_adding(schedule, adav, tcv,
loss)}~; \texttt{breakeven\_natural\_volume(current, next, threshold, loss)};
\texttt{all\_in\_per\_share(\ldots)}.
\bfield{Règles} Taux par action, signés~: un taux positif est un coût. Le meilleur
palier qualifié s'applique à toutes les actions ajoutées de la période. Les seuils sont des
parts du volume consolidé total, si bien que le moteur prend ce volume en
entrée et une prévision de celui-ci pour le mois courant. Les barèmes sont des données avec une
date d'effet, jamais du code~: ils changent chaque mois.
\bfield{Tests d'acceptation} \texttt{code/firm/feesched/tests/}~: recherche du palier
aux seuils et autour~; la rupture à 0.20~\%~; le remplissage qui fait perdre de l'argent et le
remplissage qui rapporte, avec le seuil de rentabilité en forme fermée~; coûts tout compris.
\bfield{Pour aller plus loin} Analyser la table des paliers d'un barème réel à partir de sa
page publiée vers le format de données, avec un différentiel par rapport au mois précédent.
\end{build}

\begin{omsources}
\item Cboe BZX Equities, \emph{Fee Schedule}, en vigueur à compter du 1\up{er}~septembre 2026.
\item 17 CFR 240.15c3-1 (exigences de capitaux nets pour les courtiers ou dealers).
\item New York Stock Exchange, Rule 107B (fournisseurs de liquidité supplémentaires)
tel que décrit dans les dépôts de frais de la bourse auprès de la SEC.
\end{omsources}

\section{Exercices}

\begin{exercise}[$\star$]\label{exo:m1:getting-access-equities:1}
Un membre ajoute 9 millions d'actions par jour lorsque le volume consolidé total est de 12
milliards. Trouver son palier selon l'\cref{dat:m1:getting-access-equities:bzx}, sa
ristourne quotidienne et sa ristourne pour un mois de 21 jours.
\end{exercise}

\begin{exercise}[$\star$]\label{exo:m1:getting-access-equities:2}
Calculer la ristourne quotidienne pour 23.9 et pour 24.0 millions d'actions ajoutées. Quelle est la
valeur de la cent-millième action~?
\end{exercise}

\begin{exercise}[$\star$]\label{exo:m1:getting-access-equities:3}
Avec une compensation à 2 cents et des frais réglementaires de 1 cent par tranche de 100 actions, donner le
coût tout compris pour 100 actions d'une exécution passive et d'une exécution agressive
pour le membre de l'exercice 1, et d'un aller-retour qui entre passivement et
sort agressivement.
\end{exercise}

\begin{exercise}[$\star\star$]\label{exo:m1:getting-access-equities:4}
Une société ne négocie que pour son propre compte et fait office de \omterm{def:m1:what-a-trading-firm-does:market-maker}{teneur de marché} sur 300 titres
cotés à plus de 5~\$. En utilisant l'\cref{dat:m1:getting-access-equities:netcap}, donner son
capital net minimum. Avec 2~millions de dollars de capital net, quel est son
endettement agrégé maximal~? Pourquoi aucun de ces deux chiffres n'est celui qui limite son
activité de négociation~?
\end{exercise}

\begin{exercise}[$\star\star$]\label{exo:m1:getting-access-equities:5}
Un \omterm{def:m1:getting-access-equities:dmm}{fournisseur de liquidité supplémentaire} doit se tenir au meilleur cours acheteur ou vendeur pendant 10~\% de la journée pour
chaque titre qui lui est assigné. Combien de minutes cela représente-t-il~? Une société se tient au meilleur
cours acheteur pendant 6~\% de la journée et au meilleur cours vendeur pendant 7~\%, jamais les deux à la fois. Remplit-elle
les conditions~? Que se passe-t-il si les deux coïncident toujours~?
\end{exercise}

\begin{exercise}[$\star\star$]\label{exo:m1:getting-access-equities:6}
Sur la \omterm{def:m1:us-equity-market-structure:makertaker}{plateforme inversée} de la \cref{fig:m1:getting-access-equities:venues}, un
ordre au repos coûte 13 cents pour 100 actions, tous frais compris, contre
un revenu de 17 sur la plateforme maker-taker. Donner deux raisons pour lesquelles
un \omterm{def:m1:what-a-trading-firm-does:market-maker}{teneur de marché} pourrait malgré tout laisser des ordres au repos sur cette plateforme.
\end{exercise}

\begin{exercise}[$\star\star\star$]\label{exo:m1:getting-access-equities:7}
\emph{Programmation.} Avec \texttt{breakeven\_natural\_volume}, déterminer le volume naturel
au-delà duquel le palier à 0.20~\% devient rentable, pour une perte de 27 cents et de 25
cents pour 100 actions ajoutées. Vérifier chacun avec
\texttt{extra\_volume\_worth\_adding} juste au-dessus et juste en dessous.
\end{exercise}

\begin{exercise}[$\star\star\star$]\label{exo:m1:getting-access-equities:8}
\emph{Trouver la faille.} Un plan d'affaires indique~: «~Notre stratégie passive rapporte 10 cents
pour 100 actions avant frais. La bourse verse 31 cents pour 100 actions
au titre de la liquidité ajoutée. Net~: 41 cents pour 100 actions, sur 2 millions d'actions par jour.~»
Repérer deux erreurs et recalculer.
\end{exercise}

\section{Problème~: atteindre le palier}

\begin{problem}[{Problème du week-end --- la dernière semaine du mois}]\label{pb:m1:getting-access-equities:1}
Un membre ajoute naturellement 20~millions d'actions par jour selon le barème
de l'\cref{dat:m1:getting-access-equities:bzx}~; le volume consolidé total est de 12
milliards par jour~; le mois compte 21 jours. Un volume passif supplémentaire peut être généré
en cotant plus agressivement, au prix d'une perte de 27 cents pour 100 actions avant
rabais.

\medskip\noindent\textbf{Partie I --- L'effet de seuil.}
\begin{enumerate}
\item Donner le palier du membre et son rabais mensuel.
\item Donner le volume nécessaire pour atteindre le palier suivant et le volume supplémentaire par jour.
\item Donner le gain ou la perte quotidienne liée à l'ajout de volume jusqu'à ce palier.
\item Donner le volume naturel au-delà duquel cet ajout de volume devient rentable.
\item Même question pour le palier à 0.25~\%, en partant de 24~millions. Un ajout de
volume de 24 à 30~millions est-il rentable~?
\end{enumerate}

\medskip\noindent\textbf{Partie II --- Membre ou client~?}
Un courtier dont les clients atteignent ensemble le palier à 0.25~\% propose de
reverser 80~\% de son rabais. L'adhésion coûte à la société 25\,000~\$ par mois en
frais fixes (adhésion, ports, connectivité, conformité).
\begin{enumerate}[resume]
\item Donner le rabais effectif par action pour le client.
\item À 20~millions d'actions par jour, comparer les deux voies par mois.
\item À 30~millions, refaire la comparaison.
\item À quoi d'autre la société renonce-t-elle en tant que client~?
\item À quoi d'autre s'engage-t-elle en tant que membre~?
\end{enumerate}

\medskip\noindent\textbf{Partie III --- Le mois en cours.}
\begin{enumerate}[resume]
\item Le 15 du mois, la société affiche une moyenne de 23~millions d'actions par jour. Quelle
moyenne lui faut-il sur les six derniers jours pour atteindre 24~millions sur le mois~?
\item Le volume consolidé total s'élève finalement à 13~milliards, et non 12. Quel est
le seuil désormais, et quel effet cela a-t-il sur le plan~?
\item Comment le backtester doit-il imputer les frais à une stratégie dont le volume
détermine le palier de la société~?
\item Pourquoi les bourses conçoivent-elles des barèmes avec des effets de seuil~?
\end{enumerate}

\medskip\noindent\textbf{Partie IV --- Jugement.}
\begin{enumerate}[resume]
\item Un volume échangé uniquement pour atteindre un palier nuit-il à quelqu'un~?
\item Un concurrent dix fois plus grand atteint le palier supérieur. Quantifier son
avantage par action et dire où il se manifeste sur le marché.
\item Qu'impliquerait pour ce chapitre un plafonnement des rabais par le régulateur~?
\item Lister ce que la société doit mesurer avant de négocier avec un courtier.
\item Énoncer le \emph{résultat nommé}~: le volume naturel au-delà duquel le meilleur
palier se rentabilise.
\item En une phrase~: que récompense un barème de frais~?
\end{enumerate}
\end{problem}

\section{Questions d'entretien}

\begin{interviewq}[$\star$ \iqroles{trader, développeur, chercheur}]\label{iq:m1:getting-access-equities:1}
Quelle est la différence entre un \omterm{def:m1:getting-access-equities:brokers}{courtier exécutant} et un \omterm{def:m1:getting-access-equities:brokers}{courtier compensateur}~?
\end{interviewq}

\begin{interviewq}[$\star$ \iqroles{trader, chercheur}]\label{iq:m1:getting-access-equities:2}
Qu'est-ce qu'un barème de frais maker-taker, et qu'est-ce qu'un \omterm{def:m1:getting-access-equities:tier}{palier de volume}~?
\end{interviewq}

\begin{interviewq}[$\star\star$ \iqroles{trader, chercheur}]\label{iq:m1:getting-access-equities:3}
Pourquoi quelqu'un laisserait-il un ordre au repos sur une plateforme qui facture
l'ajout de liquidité~?
\end{interviewq}

\begin{interviewq}[$\star\star$ \iqroles{chercheur, ingénieur ML}]\label{iq:m1:getting-access-equities:4}
Comment les frais des bourses doivent-ils être intégrés dans un backtest~?
\end{interviewq}

\begin{interviewq}[$\star\star$ \iqroles{trader, banque}]\label{iq:m1:getting-access-equities:5}
Vous lancez une \omterm{def:m1:what-a-trading-firm-does:prop}{société de trading pour compte propre} sur les actions américaines. Décrivez-moi
les décisions relatives à l'accès.
\end{interviewq}

\begin{interviewq}[$\star\star\star$ \iqroles{trader, chercheur}]\label{iq:m1:getting-access-equities:6}
Nous sommes le 25 du mois et il vous manque 4~\% pour atteindre un palier. Que faites-vous,
et que ne faites-vous pas~?
\end{interviewq}
